% GENERATED FILE. Edit canonical lessons, then run npm run generate:books.
% canonical-source-hash: fnv1a64:0bd0b59169c62a91
% canonical-lessons: TA-C314-tarunam, TA-C314-ninaivukur, TA-C314-nalataivil, TA-C314-atarkul, TA-C314-itaiyatuttu

\chapter{A Moment to Remember}
\label{ch:ta-a-moment-to-remember}
\hlchaptermodality{314}

\begin{chapteropening}
\textbf{By the end of this chapter:} \emph{I can talk about moments I remember.}

Complete the last lesson of chapter 314: I can talk about moments I remember.
\end{chapteropening}

\section[\texorpdfstring{\ta{தருணம்} (taruṇam)}{தருணம் (taruṇam)}]{\ta{தருணம்} (taruṇam) — a moment; an occasion}

\label{lesson:TA-C314-tarunam}



\begin{quote}
Before the new one: say the Tamil for at long last, then the Tamil for earlier.
\end{quote}

\subsection*{You'll want to know: \ta{தருணம்}}
\textbf{\ta{தருணம்}} — \emph{taruṇam} — \textquotedblleft{}a moment; an occasion\textquotedblright{}.

\begin{quote}
\ta{அது} \ta{என்} \ta{வாழ்க்கையின்} \ta{மிக} \ta{மகிழ்ச்சியான} \ta{தருணம்}.
\end{quote}

\begin{quote}
\emph{atu eṉ vāḻkkaiyiṉ mika makiḻcciyāṉa taruṇam.}
\end{quote}

\begin{quote}
That was the happiest moment of my life.
\end{quote}

Also the right moment: \ta{சரியான} \ta{தருணம்}.

\begin{cousinweb}
\textbf{\ta{தருணம்}} (\emph{taruṇam}), \textquotedblleft{}a moment, an occasion\textquotedblright{}, is a learned Tamil word whose origin is not settled in the sources at hand.
\end{cousinweb}

\begin{grammarlens}[title={What you've built}]
One more word for telling what happened.
\end{grammarlens}

\subsection*{Your turn}
\begin{itemize}
\raggedright
  \item \emph{Say it:} \emph{taruṇam}
  \item \emph{Say it:} \emph{taruṇam}, once more
  \item \emph{Say it:} \emph{atu eṉ vāḻkkaiyiṉ mika makiḻcciyāṉa taruṇam.} — That was the happiest moment of my life.
  \item \emph{Recall:} say the Tamil for at long last, then the Tamil for earlier, then say \emph{taruṇam} again
\end{itemize}

\begin{tcolorbox}[breakable,colback=teal!4,colframe=teal!35!black,title={Before you move on}]
What does \ta{தருணம்} mean? (\textquotedblleft{}A moment; an occasion\textquotedblright{}.) Say it once more.
\end{tcolorbox}
\section[\texorpdfstring{\ta{நினைவுகூர்} (niṉaivukūr)}{நினைவுகூர் (niṉaivukūr)}]{\ta{நினைவுகூர்} (niṉaivukūr) — to recall; to remember}

\label{lesson:TA-C314-ninaivukur}



\begin{quote}
Before the new one: say the Tamil for earlier, then the Tamil for a moment.
\end{quote}

\subsection*{You'll want to know: \ta{நினைவுகூர்}}
\textbf{\ta{நினைவுகூர்}} — \emph{niṉaivukūr} — \textquotedblleft{}to recall; to remember\textquotedblright{}.

\begin{quote}
\ta{தாத்தா} \ta{தன்} \ta{பள்ளி} \ta{நாட்களை} \ta{நினைவுகூர்ந்தார்}.
\end{quote}

\begin{quote}
\emph{tāttā taṉ paḷḷi nāṭkaḷai niṉaivukūrntār.}
\end{quote}

\begin{quote}
Grandfather recalled his school days.
\end{quote}

Past (I …): \ta{நினைவுகூர்ந்தேன்}, \emph{niṉaivukūrntēṉ}.

A written word; in speech \ta{நினைத்துப்} \ta{பார்}, to think back.

\begin{cousinweb}
\textbf{\ta{நினைவுகூர்}} (\emph{niṉaivukūr}) is a Tamil compound built from \ta{நினைவு}, \textquotedblleft{}remembrance\textquotedblright{} (DEDR 3683), and \ta{கூர்}: DEDR has \ta{கூர்}, \textquotedblleft{}to be sharp, keen\textquotedblright{} (DEDR 1898), and \ta{கூர்}, \textquotedblleft{}to abound\textquotedblright{} (DEDR 1899), and which of them is in this compound is not settled.
\end{cousinweb}

\begin{grammarlens}[title={What you've built}]
One more word for telling what happened.
\end{grammarlens}

\subsection*{Your turn}
\begin{itemize}
\raggedright
  \item \emph{Say it:} \emph{niṉaivukūr}
  \item \emph{Say it:} \emph{niṉaivukūr}, once more
  \item \emph{Say it:} \emph{tāttā taṉ paḷḷi nāṭkaḷai niṉaivukūrntār.} — Grandfather recalled his school days.
  \item \emph{Recall:} say the Tamil for earlier, then the Tamil for a moment, then say \emph{niṉaivukūr} again
\end{itemize}

\begin{tcolorbox}[breakable,colback=teal!4,colframe=teal!35!black,title={Before you move on}]
What does \ta{நினைவுகூர்} mean? (\textquotedblleft{}To recall; to remember\textquotedblright{}.) Say it once more.
\end{tcolorbox}
\section[\texorpdfstring{\ta{நாளடைவில்} (nāḷaṭaivil)}{நாளடைவில் (nāḷaṭaivil)}]{\ta{நாளடைவில்} (nāḷaṭaivil) — over time}

\label{lesson:TA-C314-nalataivil}



\begin{quote}
Before the new one: say the Tamil for a moment, then the Tamil for to recall.
\end{quote}

\subsection*{You'll want to know: \ta{நாளடைவில்}}
\textbf{\ta{நாளடைவில்}} — \emph{nāḷaṭaivil} — \textquotedblleft{}over time\textquotedblright{}.

\begin{quote}
\ta{நாளடைவில்} \ta{அவன்} \ta{புதிய} \ta{ஊருக்குப்} \ta{பழகினான்}.
\end{quote}

\begin{quote}
\emph{nāḷaṭaivil avaṉ putiya ūrukkup paḻakiṉāṉ.}
\end{quote}

\begin{quote}
Over time he got used to the new town.
\end{quote}

\ta{நாள்}, day + \ta{அடைவு}, reaching + -\ta{இல்}: as the days go by.

\begin{cousinweb}
\textbf{\ta{நாளடைவில்}} (\emph{nāḷaṭaivil}) is built from native \ta{நாள்}, \textquotedblleft{}day\textquotedblright{} (DEDR 3656), and \ta{அடைவு}, \textquotedblleft{}reaching\textquotedblright{}, from \ta{அடை}, \textquotedblleft{}to reach\textquotedblright{} (DEDR 79), with -\ta{இல்}.
\end{cousinweb}

\begin{grammarlens}[title={What you've built}]
One more word for telling what happened.
\end{grammarlens}

\subsection*{Your turn}
\begin{itemize}
\raggedright
  \item \emph{Say it:} \emph{nāḷaṭaivil}
  \item \emph{Say it:} \emph{nāḷaṭaivil}, once more
  \item \emph{Say it:} \emph{nāḷaṭaivil avaṉ putiya ūrukkup paḻakiṉāṉ.} — Over time he got used to the new town.
  \item \emph{Recall:} say the Tamil for a moment, then the Tamil for to recall, then say \emph{nāḷaṭaivil} again
\end{itemize}

\begin{tcolorbox}[breakable,colback=teal!4,colframe=teal!35!black,title={Before you move on}]
What does \ta{நாளடைவில்} mean? (\textquotedblleft{}Over time\textquotedblright{}.) Say it once more.
\end{tcolorbox}
\section[\texorpdfstring{\ta{அதற்குள்} (ataṟkuḷ)}{அதற்குள் (ataṟkuḷ)}]{\ta{அதற்குள்} (ataṟkuḷ) — by then; by that time}

\label{lesson:TA-C314-atarkul}



\begin{quote}
Before the new one: say the Tamil for to recall, then the Tamil for over time.
\end{quote}

\subsection*{You'll want to know: \ta{அதற்குள்}}
\textbf{\ta{அதற்குள்}} — \emph{ataṟkuḷ} — \textquotedblleft{}by then; by that time\textquotedblright{}.

\begin{quote}
\ta{நான்} \ta{நிலையத்துக்கு} \ta{வந்தேன்}, \ta{ஆனால்} \ta{அதற்குள்} \ta{ரயில்} \ta{போய்விட்டது}.
\end{quote}

\begin{quote}
\emph{nāṉ nilaiyattukku vantēṉ, āṉāl ataṟkuḷ rayil pōyviṭṭatu.}
\end{quote}

\begin{quote}
I got to the station, but by then the train had gone.
\end{quote}

\ta{அது} + -\ta{க்கு} + \ta{உள்}, within that; after a verb, \ta{வருவதற்குள்}, before … could come.

\begin{cousinweb}
\textbf{\ta{அதற்குள்}} (\emph{ataṟkuḷ}) is built from \ta{அது}, \textquotedblleft{}that\textquotedblright{} (DEDR 1), and \ta{உள்}, \textquotedblleft{}inside\textquotedblright{} (DEDR 698): within that time.
\end{cousinweb}

\begin{grammarlens}[title={What you've built}]
One more word for telling what happened.
\end{grammarlens}

\subsection*{Your turn}
\begin{itemize}
\raggedright
  \item \emph{Say it:} \emph{ataṟkuḷ}
  \item \emph{Say it:} \emph{ataṟkuḷ}, once more
  \item \emph{Say it:} \emph{nāṉ nilaiyattukku vantēṉ, āṉāl ataṟkuḷ rayil pōyviṭṭatu.} — I got to the station, but by then the train had gone.
  \item \emph{Recall:} say the Tamil for to recall, then the Tamil for over time, then say \emph{ataṟkuḷ} again
\end{itemize}

\begin{tcolorbox}[breakable,colback=teal!4,colframe=teal!35!black,title={Before you move on}]
What does \ta{அதற்குள்} mean? (\textquotedblleft{}By then; by that time\textquotedblright{}.) Say it once more.
\end{tcolorbox}
\section[\texorpdfstring{\ta{இதையடுத்து} (itaiyaṭuttu)}{இதையடுத்து (itaiyaṭuttu)}]{\ta{இதையடுத்து} (itaiyaṭuttu) — following this}

\label{lesson:TA-C314-itaiyatuttu}



\begin{quote}
Before the new one: say the Tamil for over time, then the Tamil for by then.
\end{quote}

\subsection*{You'll want to know: \ta{இதையடுத்து}}
\textbf{\ta{இதையடுத்து}} — \emph{itaiyaṭuttu} — \textquotedblleft{}following this\textquotedblright{}.

\begin{quote}
\ta{இதையடுத்து}, \ta{காவல்துறை} \ta{புலனாய்வைத்} \ta{தொடங்கியது}.
\end{quote}

\begin{quote}
\emph{itaiyaṭuttu, kāvaltuṟai pulaṉāyvait toṭaṅkiyatu.}
\end{quote}

\begin{quote}
Following this, the police began an investigation.
\end{quote}

A news-style link: \ta{இதை} + \ta{அடுத்து}, next to this.

\begin{cousinweb}
\textbf{\ta{இதையடுத்து}} (\emph{itaiyaṭuttu}) is built from \ta{இதை}, \textquotedblleft{}this\textquotedblright{}, and \ta{அடுத்து}, \textquotedblleft{}next to\textquotedblright{}, from native \ta{அடு}, \textquotedblleft{}to be next, near\textquotedblright{} (DEDR 79).
\end{cousinweb}

\begin{grammarlens}[title={What you've built}]
One more word for telling what happened.
\end{grammarlens}

\subsection*{Your turn}
\begin{itemize}
\raggedright
  \item \emph{Say it:} \emph{itaiyaṭuttu}
  \item \emph{Say it:} \emph{itaiyaṭuttu}, once more
  \item \emph{Say it:} \emph{itaiyaṭuttu, kāvaltuṟai pulaṉāyvait toṭaṅkiyatu.} — Following this, the police began an investigation.
  \item \emph{Recall:} say the Tamil for over time, then the Tamil for by then, then say \emph{itaiyaṭuttu} again
\end{itemize}

\begin{tcolorbox}[breakable,colback=teal!4,colframe=teal!35!black,title={Before you move on}]
What does \ta{இதையடுத்து} mean? (\textquotedblleft{}Following this\textquotedblright{}.) Say it once more.
\end{tcolorbox}
